\subsection{Controle de atitude da espaçonave}
\label{sec_controle_PD_atitude}

A atitude da espaçonave é descrita por um quaternion unitário, conforme a Equação~\eqref{eq_qrel_sincronizacao}, ou seja:

\begin{equation}
\vec q_B =
\begin{bmatrix}
q_{B,w} \\
q_{B,x} \\
q_{B,y} \\
q_{B,z}
\end{bmatrix}.
\end{equation}

Por sua vez, a referência de atitude é a atitude da espaçonave alvo, sendo especificada por um quaternion desejado, como:

\begin{equation}
\vec q_{B,\text{ref}} =
\begin{bmatrix}
q_{B,w,\text{ref}} \\
q_{B,x,\text{ref}} \\
q_{B,y,\text{ref}} \\
q_{B,z,\text{ref}}
\end{bmatrix}.
\end{equation}

O erro de atitude é obtido a partir do quaternion relativo, isto é:

\begin{equation}
\vec q_e
=
\vec q_{B,\text{ref}} \otimes \vec q_B^{-1}.
\label{eq_controle_PD_quat}
\end{equation}

Adota-se a convenção $q_{e,w}\ge 0$. Caso $q_{e,w}<0$, define-se $\vec q_e \leftarrow -\vec q_e$.
Definindo a parte vetorial do quaternion de erro como
$\vec q_{e,v} = [\,q_{e,x}\;\;q_{e,y}\;\;q_{e,z}\,]^T$,
a lei de controle \gls{pd} de atitude atua sobre $\vec q_{e,v}$ e sobre a velocidade angular:

\begin{equation}
\vec\tau_B
=
-K_{p,B}\,\vec q_{e,v}
-
K_{d,B}\,\vec\omega_B.
\label{eq_controle_PD_atitude}
\end{equation}

Em que \(K_{p,B}\) e \(K_{d,B}\) são matrizes de ganhos proporcionais e derivativos, respectivamente.
Por sua vez, a velocidade angular da espaçonave, expressa em $\{B\}$, é escrita como:

\begin{equation}
\vec\omega_B =
\begin{bmatrix}
\omega_x \\
\omega_y \\
\omega_z
\end{bmatrix}.
\end{equation}

A ação de controle $\vec\tau_B$ é inserida no vetor de torques generalizados $\vec\tau$ das equações de movimento presentes na Seção~\ref{sec_atitude_inercial}, especificamente na Equação~\eqref{eq_dinamica_geral}, já que as primeiras coordenadas generalizadas estão associadas à atitude da espaçonave.
Dessa forma, o controlador \gls{pd} opera diretamente sobre a parcela correspondente aos torques da espaçonave no vetor $\vec\tau$.
Reorganizando os termos associados às coordenadas do \gls{mr} na
Equação~\eqref{eq_dinAcoplada_espaconave_compacta}, obtém-se a dinâmica da espaçonave explicitando o efeito de reação do manipulador:

\begin{equation}
M_{BB}(\vec y\,)\,\ddot{\vec\eta}_B
+
C_{BB}(\vec y\,,\dot{\vec y\,})\,\dot{\vec\eta}_B
=
\vec\tau_B
-
\underbrace{\bigl[
M_{B\xi}(\vec y\,)\,\ddot{\vec\xi}
+
C_{B\xi}(\vec y\,,\dot{\vec y\,})\,\dot{\vec\xi}
\bigr]}_{\displaystyle \vec\tau_{\text{reac}}}.
\label{eq_dina_espaconave_reac_compacta}
\end{equation}


Em que:

\begin{equation}
\vec\tau_{\text{reac}}
=
M_{B\xi}(\vec y\,)\,\ddot{\vec\xi}
+
C_{B\xi}(\vec y\,,\dot{\vec y\,})\,\dot{\vec\xi}.
\label{eq_def_tau_reac}
\end{equation}


A Equação~\ref{eq_def_tau_reac} representa o torque de reação na espaçonave devido ao movimento das juntas do manipulador.
Durante o momento em que o \gls{mr} está no modo não operacional:

\begin{equation}
\dot{\vec\xi} = \vec 0,
\qquad
\ddot{\vec\xi} = \vec 0.
\end{equation}


Tem-se $\vec\tau_{\text{reac}} = \vec 0$.
Quando o \gls{mr} se movimenta, os termos de acoplamento $M_{B\xi}$ e $C_{B\xi}$ introduzem $\vec\tau_{\text{reac}}$ como perturbação interna, de modo que o controlador \gls{pd} em \eqref{eq_controle_PD_atitude} deve simultaneamente reduzir o erro de atitude e compensar os torques de reação induzidos pelo manipulador, preservando o alinhamento da espaçonave requerido para a operação de ancoramento.
